\section{Experimental Setup: Source-Native Replay}
\label{sec:evidence}
TimeSage-EV motivates the three recurring mechanism families---cutoff validity, release alignment, and exogenous grounding---and exposes persistent skills to later periods~\citep{yao2026timesage}. Our experiment is not an exact rerun of its unavailable evaluated snapshot. Instead, we reconstruct content-addressed endpoints directly from first-party institutional release histories and intervene on the same reusable operations.

\subsection{Frozen endpoints and scorers}
The initial replay contains 23 endpoints from BEA GDP releases and NOAA-CPC ENSO Diagnostic Discussions~\citep{bea2026archive,noaa2026enso}: twelve pre-April endpoints form C3-R and eleven April--May endpoints form held-out C4-R. Raw pages are SHA-256 bound, evidence identifiers are opaque, and grounding exposes the full content-region sentence pool rather than a gold-only candidate set. Family scorers are deterministic and frozen before model outcomes; a pre-execution self-test gives 23/23 success on gold outputs and 0/23 on constructed negatives.

A post-review validation adds 12 EIA Weekly Petroleum Status Report cutoff endpoints from a third institutional system~\citep{eia2026wpsrarchive}. EIA is chosen by a structural compatibility rule---recurring first-party releases, explicit publication dates, represented data-ending weeks, and compatibility with the already frozen cutoff interface---before any EIA model outcome is opened. Four endpoints form an earlier diagnostic split and eight April--May endpoints form held-out validation. T1 and G1 are reused byte-for-byte; source manifest, endpoints, scorer, harness, and condition order are frozen before execution. This is informed post-review validation, not an independently preregistered confirmation.

\subsection{Questions and inference contract}
RQ1 asks whether adding the no-skill anchor changes a verdict that would look positive under targeted-versus-generic comparison alone. RQ2 asks whether targeted procedures retain a residual over behaviorally neutral extra computation after G$_0$ is independently validated against fresh N. RQ3 asks whether byte-identical procedures retain direction on compatible held-out releases while preserving ceiling and incompatible-domain nulls. RQ4 asks whether callable skill state adds material value over an operation-matched retrieval/context treatment R that exposes exactly the same frozen temporal-operation output.

The endpoint is the independent scientific unit. Repeats characterize execution robustness and are averaged within endpoint; they are never counted as additional tasks. For non-ceiling cells we report endpoint win/tie/loss patterns, percentile bootstrap intervals over endpoints, and exact two-sided sign tests over non-ties when resolution permits. Two- and three-endpoint cells are explicitly descriptive. We do not pool heterogeneous procedures, models, or domains into an omnibus procedure-effect statistic.

\subsection{Evidence layers and parity}
DeepSeek-V4-Pro is the registered primary track with five-repeat historical replay plus fresh identification runs. A pre-outcome support rule selects Kimi-K3 as a separately labeled replication; a later hash audit removes one overwritten historical repeat block without replacement. Doubao-Seed-2.0-Mini is an exploratory capacity stress. The EIA layer evaluates DeepSeek and the mini model. After adding the prospectively frozen G$_0$ and R controls, the retained source-native evidence ledger contains 2,056 runtime-valid rows over the same 35 distinct endpoints from BEA, NOAA, and EIA. The endpoint count, not the enlarged repeat/control-row count, determines scientific resolution.

Every targeted procedure implementation has a family-matched generic helper with the same callable signature, identical Python lexical-token count and AST-node count, and equal external tool-call budget. G$_0$ uses the same helper-execution surface but is information-empty: it returns \texttt{\{\}} and reads neither package nor context. For R, zero-call parity verifies on every tested endpoint that removing the added retrieval-operation field recovers the original evidence package exactly and that the materialized operation output is content-identical to T's helper output. Full execution chronology, content hashes, per-call CSV checkpoints, and provenance bindings are reported in Appendix~A and the supplement.
